\documentclass{article}
\usepackage[utf8]{inputenc}
\usepackage{geometry}
 \geometry{
 a4paper,
 total={170mm,257mm},
 left=20mm,
 top=20mm,
 }
 \usepackage{graphicx}
 \usepackage{titling}
 \usepackage{booktabs}
 \usepackage{listings}
 \usepackage{xcolor}
 \usepackage{hyperref}
 \usepackage{fancyvrb}
 \usepackage{float}
 \usepackage{amsmath}
 \definecolor{darkblue}{rgb}{0.0, 0.0, 0.5}

\RecustomVerbatimEnvironment{verbatim}{Verbatim}{formatcom=\color{darkblue}}

\setlength{\parindent}{0pt} %lui dau dong


\title{\textbf{DEEP LEARNING} \\ 
{\Large \textbf{Labwork 3 - Logistic Regression}}}

\date{May 2026}
 
 \usepackage{fancyhdr}
\fancypagestyle{plain}{
    \fancyhf{} 
    \fancyfoot[L]{\thedate}
    \fancyfoot[R]{\thepage}
    \fancyhead[L]{University of Science and Technology of Hanoi}
    \fancyhead[R]{Deep Learning}
}
\makeatletter
\def\@maketitle{%
  \newpage
  % \null
  \vskip 1em %
  \begin{center}%
  \let \footnote \thanks
    {\LARGE \@title \par}%
    \vskip 1em%
    %{\large \@date}%
  \end{center}%
  \par
  \vskip 1em}
\makeatother

\usepackage{cmbright}

% Configuration for code blocks
\definecolor{codegreen}{rgb}{0,0.6,0}
\definecolor{backcolour}{rgb}{0.95,0.95,0.92}
\lstset{
    backgroundcolor=\color{backcolour},   
    commentstyle=\color{codegreen},
    basicstyle=\ttfamily\footnotesize,
    breaklines=true,
    showstringspaces=false
}

\begin{document}

\pagestyle{plain}
\maketitle

\noindent\begin{tabular}{@{}ll}
    Student: & Foo Cheung Nicolas - ES2540023 \\
\end{tabular}

\section{Algorithm Implementation}

In this lab, we implemented logistic regression from scratch to predict whether a person would get a loan based on their experience and salary. The goal is to understand the effect of gradient descent optimization and the role of the learning rate in convergence.



\subsection{Model Definition}

The logistic regression model predicts the probability of a positive outcome using the sigmoid function:
\[
\hat{y} = \sigma(z) = \frac{1}{1 + e^{-z}}, \quad z = w_1 x_1 + w_2 x_2 + w_0
\]
\subsection{Loss Function}

The binary cross-entropy loss for $N$ data points is:
\[
J(w_0, w_1, w_2) = \frac{1}{N} \sum_{i=1}^N L_i = -\frac{1}{N} \sum_{i=1}^N \Big[ y_i \log(\hat{y}_i) + (1-y_i) \log(1-\hat{y}_i) \Big]
\]

\subsection{Gradients}
The partial derivatives of the loss function with respect to the weights are:
\[
\frac{\partial J}{\partial w_0} = \frac{1}{N} \sum_{i=1}^N (\hat{y}_i - y_i)
\]
\[
\frac{\partial J}{\partial w_1} = \frac{1}{N} \sum_{i=1}^N (\hat{y}_i - y_i) x_i^{(1)}
\]
\[
\frac{\partial J}{\partial w_2} = \frac{1}{N} \sum_{i=1}^N (\hat{y}_i - y_i) x_i^{(2)}
\]

\subsection{Gradient Descent}

We update the weights iteratively:
\[
w_j \gets w_j - r \frac{\partial J}{\partial w_j}, \quad j = 0,1,2
\]
where $r$ is the learning rate. 

\section{Effect of Learning Rate}

The learning rate $r$ significantly affects convergence:

\begin{itemize}
    \item \textbf{Small learning rate ($r = 0.001$):}
          Very slow convergence, many iterations needed.

    \item \textbf{Moderate learning rate ($r = 0.05$):}
          Fast and stable convergence.

    \item \textbf{Large learning rate ($r = 0.1$):}
          May overshoot minimum, loss oscillates, may diverge.
\end{itemize}

Choosing an appropriate learning rate is crucial for gradient descent efficiency and stability.

\section{Result Visualization}

The graph shows the logistic regression decision boundary for predicting loan approval based on Experience and Salary. Red points indicate profiles approved for a loan, blue points indicate profiles refused, and the green line represents the decision boundary learned by gradient descent.


\begin{figure}[H]
        \centering
        \includegraphics[width=0.75\linewidth]{logistic.png}
        \caption{Logistic Regression graph}
        \label{fig:placeholder}
\end{figure}

\end{document}